% Source PDF pages 65--66; manual pages 2-36--2-37.
\section{Equations of Motion}\label{sec:equations-of-motion}

The equations of motion are derived from Newton's second law applied to a point
mass:
\begin{equation}
  \vec a_I=\frac{\sum\vec F}{m}.
  \label{eq:newton-point-mass}
\end{equation}
This equation defines acceleration in an inertial coordinate system. For trajectory
analysis it is more convenient to express acceleration in the earth-fixed system,
$\vec a_{\oplus}$. The earth-fixed system rotates relative to the inertial system at
the earth spin rate ${}_{I}\vec\omega_{\oplus}$.

In inertial coordinates, acceleration is the second derivative of the vehicle position
vector:
\begin{equation}
  \vec a_I
  ={}^{I}\!\frac{d^2\vec r}{dt^2}
  ={}^{I}\!\frac{d}{dt}
    \left({}^{I}\!\frac{d\vec r}{dt}\right).
  \label{eq:inertial-position-second-derivative}
\end{equation}
Here ${}^{I}d/dt$ denotes differentiation with respect to inertial coordinates. If
the inner derivative is evaluated relative to the rotating earth-fixed frame, the
transport theorem gives
\begin{equation}
  \vec a_I
  ={}^{I}\!\frac{d}{dt}
    \left[
      {}^{\oplus}\!\frac{d\vec r}{dt}
      +\left({}_{I}\vec\omega_{\oplus}\times\vec r\right)
    \right].
  \label{eq:inertial-derivative-earth-fixed-inner}
\end{equation}
The notation ${}^{\oplus}d/dt$ denotes differentiation with respect to earth-fixed
coordinates. Differentiating the outer derivative in the same rotating frame gives
\begin{equation}
  \begin{aligned}
  \vec a_I={}&
  {}^{\oplus}\!\frac{d}{dt}
    \left[
      {}^{\oplus}\!\frac{d\vec r}{dt}
      +{}_{I}\vec\omega_{\oplus}\times\vec r
    \right]\\
  &+{}_{I}\vec\omega_{\oplus}\times
    \left[
      {}^{\oplus}\!\frac{d\vec r}{dt}
      +{}_{I}\vec\omega_{\oplus}\times\vec r
    \right].
  \end{aligned}
  \label{eq:transport-theorem-expanded}
\end{equation}
Expanding the derivatives results in
\begin{equation}
  \begin{aligned}
  \vec a_I={}&
  \left[
    {}^{\oplus}\!\frac{d^2\vec r}{dt^2}
    +\frac{d\,{}_{I}\vec\omega_{\oplus}}{dt}\times\vec r
    +{}_{I}\vec\omega_{\oplus}\times
      {}^{\oplus}\!\frac{d\vec r}{dt}
  \right]\\
  &+{}_{I}\vec\omega_{\oplus}\times
      {}^{\oplus}\!\frac{d\vec r}{dt}
   +{}_{I}\vec\omega_{\oplus}\times
      \left({}_{I}\vec\omega_{\oplus}\times\vec r\right).
  \end{aligned}
  \label{eq:inertial-acceleration-expanded}
\end{equation}
The derivative of the earth spin rate is zero because
${}_{I}\vec\omega_{\oplus}$ is constant. The notation is simplified with
\begin{equation}
  \vec V_{\oplus}
    ={}^{\oplus}\!\frac{d\vec r}{dt},
  \qquad
  \vec a_{\oplus}
    ={}^{\oplus}\!\frac{d^2\vec r}{dt^2}.
  \label{eq:earth-relative-velocity-acceleration-definitions}
\end{equation}
Substitution gives
\begin{equation}
  \begin{aligned}
  \vec a_I={}&
    \left[
      \vec a_{\oplus}
      +{}_{I}\vec\omega_{\oplus}\times\vec V_{\oplus}
    \right]\\
    &+\left[
      {}_{I}\vec\omega_{\oplus}\times\vec V_{\oplus}
      +{}_{I}\vec\omega_{\oplus}\times
        \left({}_{I}\vec\omega_{\oplus}\times\vec r\right)
    \right].
  \end{aligned}
  \label{eq:inertial-acceleration-grouped}
\end{equation}
Combining the two identical Coriolis terms produces
\begin{equation}
  \vec a_I
  =\vec a_{\oplus}
   +2\left({}_{I}\vec\omega_{\oplus}\times\vec V_{\oplus}\right)
   +{}_{I}\vec\omega_{\oplus}\times
      \left({}_{I}\vec\omega_{\oplus}\times\vec r\right).
  \label{eq:inertial-earth-acceleration-relation}
\end{equation}
Finally, combining \cref{eq:newton-point-mass} with
\cref{eq:inertial-earth-acceleration-relation} gives the earth-fixed equation of
motion:
\begin{equation}
  \vec a_{\oplus}
  =\frac{\sum\vec F}{m}
   -2\left({}_{I}\vec\omega_{\oplus}\times\vec V_{\oplus}\right)
   -{}_{I}\vec\omega_{\oplus}\times
      \left({}_{I}\vec\omega_{\oplus}\times\vec r\right).
  \label{eq:earth-fixed-equation-of-motion}
\end{equation}
This second-order differential equation can be written as the following pair of
first-order equations:
\begin{equation}
  \begin{aligned}
    \dot{\vec r}_{\oplus}&=\vec V_{\oplus},\\[0.3em]
    \dot{\vec V}_{\oplus}
      &=\frac{\sum\vec F}{m}
       -2\left({}_{I}\vec\omega_{\oplus}\times\vec V_{\oplus}\right)
       -{}_{I}\vec\omega_{\oplus}\times
          \left({}_{I}\vec\omega_{\oplus}\times\vec r\right).
  \end{aligned}
  \label{eq:first-order-equations-of-motion}
\end{equation}
Given an initial position $\vec r$, an earth-relative velocity
$\vec V_{\oplus}$, an initial time $t_0$, the vehicle mass $m$, and the forces
$\sum\vec F$, these equations can be integrated with respect to time to obtain a
trajectory. The associated integration functions are in the \texttt{derivs} module.
